% Part: sets-functions-relations
% Chapter: size-of-sets-complete

\documentclass[../../../include/open-logic-chapter]{subfiles}

\begin{document}

\olchapter{sfr}{siz}{సమితుల పరిమాణం}

\begin{editorial}
ఈ అధ్యాయం లెక్కింపులు, లెక్కించదగినతనం, లెక్కించలేనితనం గురించి చర్చిస్తుంది.
కొన్ని విభాగాలు రెండేసి రూపాల్లో ఉన్నాయి. మరింత ప్రాథమికమైన రూపం
లెక్కింపులను జాబితాలుగా, లేదా $\PosInt$ నుంచి సంగ్రస్త ప్రమేయాలుగా
తీసుకుంటుంది; మరింత అమూర్తమైన రూపం లెక్కింపులను $\Nat$తో ద్విగుణ
ప్రమేయాలుగా నిర్వచిస్తుంది.
\end{editorial}

\olimport{introduction}

\olimport{enumerability}

\olimport{zig-zag}

\olimport{pairing}

\olimport{pairing-alt}

\olimport{non-enumerability}

\olimport{reduction}

\olimport{equinumerous-sets}

\olimport{comparing-size}

\olimport{schroder-bernstein}

\begin{editorial}
కిందివైన \olref[sfr][siz][enm-alt]{sec},
\olref[sfr][siz][nen-alt]{sec}, \olref[sfr][siz][red-alt]{sec} విభాగాలు
\olref[sfr][siz][enm]{sec}, \olref[sfr][siz][nen]{sec},
\olref[sfr][siz][red]{sec}లకు ప్రత్యామ్నాయ రూపాలు. టిమ్ బటన్ తన
Open Set Theory పాఠ్యంలో వాడటానికి వీటిని రచించారు. ఇవి కొంచెం ఉన్నత
స్థాయిలో ఉండి, సమితి సిద్ధాంత సందర్భానికి మరింత అనువైన వేరొక
లెక్కించదగినతనం నిర్వచనాన్ని వాడతాయి (అంటే, జాబితా చేయగలగడం లేదా
$\PosInt$ నుంచి సంగ్రస్త ప్రమేయానికి వ్యాప్తి కావడం బదులు, $\Nat$తో
లేదా దాని ఆరంభ ఖండంతో ద్విగుణ ప్రమేయం ఉండటం).
\sourcecorrection{OLTESIZ-001}{ఆంగ్ల మూలంలోని ``difference definition''ను
సందర్భం కోరిన ``different definition''గా సరిచేసి అనువదించాం.}
\end{editorial}

\olimport{enumerability-alt}
\olimport{non-enumerability-alt}
\olimport{reduction-alt}

\OLEndChapterHook

\end{document}
